\section{Cluster Agreement Metrics}\label{sec:bg:cluster_metrics}

A clustering algorithm returns a labelling of $N$ points.
The integers used as cluster names are arbitrary: two labellings that group the same points together are the same partition.
External evaluation therefore compares a predicted partition $\hat{Y}$ with a reference partition $Y$ in a way that does not depend on those names.

All scores below are computed from the contingency table
\begin{equation}\label{eq:bg:contingency}
  n_{ij}=\sum_{n=1}^{N}\indicator[y_n=i,\,\hat{y}_n=j],
\end{equation}
with row sums $a_i=\sum_j n_{ij}$ and column sums $b_j=\sum_i n_{ij}$ equal to the sizes of the reference and predicted clusters.

\subsection*{Adjusted Rand index}

Each unordered pair of points is either together or apart in a given partition.
A cluster of size $a_i$ contributes $\binom{a_i}{2}$ co-clustered pairs, so
\begin{equation}\label{eq:bg:pair_counts}
  S_{Y\hat{Y}}=\sum_{i,j}\binom{n_{ij}}{2},\qquad
  S_Y=\sum_i\binom{a_i}{2},\qquad
  S_{\hat{Y}}=\sum_j\binom{b_j}{2}
\end{equation}
count the pairs that share a cluster in both labellings, in $Y$, and in $\hat{Y}$, out of $\binom{N}{2}$ pairs in total.

The Rand index is the fraction of pairs on which $Y$ and $\hat{Y}$ agree: together in both, or apart in both \parencite{hubert1985comparing},
\begin{equation}\label{eq:bg:rand}
  R(Y,\hat{Y})
  =
  1 - \frac{S_Y + S_{\hat{Y}} - 2\,S_{Y\hat{Y}}}{\binom{N}{2}}
  \in[0,1].
\end{equation}
The numerator $S_Y+S_{\hat{Y}}-2\,S_{Y\hat{Y}}$ is the number of pairs that are together in one partition and apart in the other.

When most points lie in different clusters, the bulk of pairs are apart in both labellings by chance, and $R$ stays close to one even for weak agreement.
The adjusted Rand index $\ARI$ subtracts the expected overlap under a permutation of labels that keeps the cluster sizes $\{a_i\}$ and $\{b_j\}$ fixed,
\begin{equation}\label{eq:bg:ari_expected}
  \E[S_{Y\hat{Y}}]
  =
  \frac{S_Y\,S_{\hat{Y}}}{\binom{N}{2}},
\end{equation}
and rescales so that exact recovery scores one \parencite{hubert1985comparing}:
\begin{equation}\label{eq:bg:ari}
  \ARI
  =
  \frac{S_{Y\hat{Y}}-\E[S_{Y\hat{Y}}]}{\tfrac{1}{2}\bigl(S_Y+S_{\hat{Y}}\bigr)-\E[S_{Y\hat{Y}}]},
  \qquad
  \ARI\le 1.
\end{equation}
Hence $\ARI=1$ at exact recovery up to relabelling, $\ARI\approx 0$ at chance, and negative values indicate worse-than-random alignment.

\subsection*{Adjusted mutual information}

Mutual information asks how much knowing one labelling reduces uncertainty about the other \parencite{vinh2010information}.
With the marginal entropies
\begin{equation}\label{eq:bg:entropy_marginals}
  H(Y)=-\sum_{i:\,a_i>0}\frac{a_i}{N}\log\frac{a_i}{N},
  \qquad
  H(\hat{Y})=-\sum_{j:\,b_j>0}\frac{b_j}{N}\log\frac{b_j}{N},
\end{equation}
the empirical mutual information is
\begin{equation}\label{eq:bg:mi}
  \mathrm{MI}(Y;\hat{Y})
  =
  \sum_{i,j:\,n_{ij}>0}\frac{n_{ij}}{N}\,
  \log\frac{N\,n_{ij}}{a_i b_j},
\end{equation}
with the convention $0\log 0=0$.
It is zero when the labellings are independent, and equals either marginal entropy when one labelling determines the other.

The expectation of $\mathrm{MI}$ under the same permutation null grows with the number of clusters.
The adjusted mutual information $\AMI$ subtracts that baseline and rescales by the average of the marginal entropies \parencite{vinh2010information},
\begin{equation}\label{eq:bg:ami}
  \AMI(Y;\hat{Y})
  =
  \frac{\mathrm{MI}(Y;\hat{Y})-\E[\mathrm{MI}]}{\tfrac{1}{2}\bigl(H(Y)+H(\hat{Y})\bigr)-\E[\mathrm{MI}]},
\end{equation}
when the denominator is positive; otherwise $\AMI\equiv 0$.
The range matches $\ARI$.
$\AMI$ weights clusters more evenly than pair counting, so large groups dominate it less.

\subsection*{Homogeneity, completeness, and V-measure}

Homogeneity is high when each predicted cluster contains points from only one reference class.
Completeness is high when each reference class is collected into a single predicted cluster.
Splitting a class across many predicted clusters preserves homogeneity and lowers completeness; merging several classes into one predicted cluster does the reverse.

With the conditional entropies
\begin{equation}\label{eq:bg:cond_ent}
  H(Y\mid \hat{Y})=-\sum_{i,j:\,n_{ij}>0}\frac{n_{ij}}{N}\log\frac{n_{ij}}{b_j},
  \qquad
  H(\hat{Y}\mid Y)=-\sum_{i,j:\,n_{ij}>0}\frac{n_{ij}}{N}\log\frac{n_{ij}}{a_i},
\end{equation}
the two quantities are \parencite{rosenberg2007v}
\begin{equation}\label{eq:bg:vmeasure_hc}
  h = 1-\frac{H(Y\mid \hat{Y})}{H(Y)},
  \qquad
  c = 1-\frac{H(\hat{Y}\mid Y)}{H(\hat{Y})},
\end{equation}
with $h=1$ if $H(Y)=0$ and $c=1$ if $H(\hat{Y})=0$.
The V-measure is their harmonic mean,
\begin{equation}\label{eq:bg:vmeasure}
  \mathrm{V}(Y,\hat{Y})=\frac{2hc}{h+c},
  \qquad h+c>0,
\end{equation}
and $\mathrm{V}=0$ if $h=c=0$.
The harmonic mean is small if either factor is small, so a high V-measure requires both.

The three families return one at exact recovery and about zero at chance.
Between those extremes, $\ARI$ follows pair agreement, $\AMI$ follows shared information, and homogeneity versus completeness separates splitting from merging.
